\paragraph{最佳几何均匀化长度上的共振保真与传输预算失配}
在黄金分支的 \(W_1\) 传输奇偶双相律与 bulk 共振梯的 Fibonacci 极限常数已经分别固定之后，还可把几何均匀化与算术共振之间的制度断裂进一步压缩为两条直接后果。其核心是：即便沿最佳逼近长度 \(N=F_n\) 观察，一维几何误差也只按 \(F_n^{-1}\) 衰减，而同一子列上的共振梯却保留严格正的宏观幅度；因此，几何意义下的准均匀化并不会迫使算术共振同步消失。进一步，对全部频率的累积共振预算 \(\sum_{u\le N}C_\varphi(u)^2\) 而言，其增长至少为对数量级，因而不可能由仅依赖一维传输差异的任何上界机制统一控制。

\begin{theorem}[最佳逼近长度上的几何均匀化与算术共振不相容]\label{thm:xi-time-part9zi-fibonacci-transport-resonance-incompatibility}
\leanverified{paper\_xi\_time\_part9zi\_fibonacci\_transport\_resonance\_incompatibility}
记
\[
T_N(\varphi^{-1}):=\mathsf T_N(\varphi^{-1})
=
W_1(\mu_N,\lambda),
\]
其中 \(\mu_N\) 为黄金 Kronecker 轨道的经验分布，\(\lambda\) 为 Lebesgue 测度。又记 \(C_\varphi(u)\) 为 bulk 共振梯。则沿 Fibonacci 长度 \(N=F_n\) 有
\[
F_n\,T_{F_n}(\varphi^{-1})
\longrightarrow
\begin{cases}
\displaystyle \frac12+\frac{1}{2\sqrt5}, & n\ \text{为偶},\\[2mm]
\displaystyle \frac12-\frac{1}{2\sqrt5}+\frac1{15}, & n\ \text{为奇},
\end{cases}
\]
并且存在严格正常数
\[
c_{\mathrm{Fib}}:=\lim_{n\to\infty}C_\varphi(F_n)>0.
\]
因此
\[
T_{F_n}(\varphi^{-1})=\Theta(F_n^{-1}),
\qquad
C_\varphi(F_n)=\Theta(1),
\]
并且
\[
\liminf_{n\to\infty}
F_n\,T_{F_n}(\varphi^{-1})\,C_\varphi(F_n)>0.
\]
故在最优一维准均匀化子序列上，算术共振并不随几何误差同步消失；最佳几何分布与共振湮灭在黄金分支上不能同时发生。
\end{theorem}
\begin{proof}
由定理 \ref{thm:xi-golden-fibonacci-w1-star-strict-parity-law}，
\[
\lim_{n\to\infty,\ n\ \mathrm{even}}
F_n\,T_{F_n}(\varphi^{-1})
=
\frac12+\frac{1}{2\sqrt5},
\]
\[
\lim_{n\to\infty,\ n\ \mathrm{odd}}
F_n\,T_{F_n}(\varphi^{-1})
=
\frac12-\frac{1}{2\sqrt5}+\frac1{15}.
\]
这两个极限常数都严格为正，故
\[
F_n\,T_{F_n}(\varphi^{-1})=\Theta(1),
\qquad
T_{F_n}(\varphi^{-1})=\Theta(F_n^{-1}).
\]

另一方面，由推论 \ref{cor:fold-resonance-ladder-fib-luc-limits}，
\[
C_\varphi(F_n)\longrightarrow c_{\mathrm{Fib}}>0,
\]
从而
\[
C_\varphi(F_n)=\Theta(1).
\]
于是
\[
\liminf_{n\to\infty}
F_n\,T_{F_n}(\varphi^{-1})\,C_\varphi(F_n)
\ge
\left(\frac12-\frac{1}{2\sqrt5}+\frac1{15}\right)c_{\mathrm{Fib}}
>0.
\]
这便得到所述不相容性。证毕。
\end{proof}

\begin{corollary}[传输正则性不能约束共振预算]\label{cor:xi-time-part9zi-transport-regularity-cannot-bound-resonance-budget}
\leanverified{paper\_xi\_time\_part9zi\_transport\_regularity\_cannot\_bound\_resonance\_budget}
定义累积共振预算
\[
\mathcal E(N):=\sum_{u=1}^{N}C_\varphi(u)^2.
\]
则存在常数 \(c>0\) 与 \(N_0\)，使得对一切 \(N\ge N_0\) 都有
\[
\mathcal E(N)\ge c\log N.
\]
与此同时，沿 Fibonacci 长度 \(N=F_n\) 有
\[
F_n\,T_{F_n}(\varphi^{-1})=O(1),
\qquad
\mathcal E(F_n)\longrightarrow\infty.
\]
特别地，不存在任何仅依赖一维传输差异 \(T_N(\varphi^{-1})\) 的上界机制，可以在最佳逼近长度上控制累积共振预算。等价地，几何意义下的准均匀采样并不驯化 bulk 共振的累积增长。
\end{corollary}
\begin{proof}
由定理 \ref{thm:xi-resonance-closure-log-lower-bound-fib-luc}，
\[
\liminf_{N\to\infty}
\frac{1}{\log N}\sum_{u=1}^{N}C_\varphi(u)^2
\ge
\frac{c_{\mathrm{Fib}}^2+c_{\mathrm{Luc}}^2}{\log\varphi}>0.
\]
因此可取某个常数 \(c>0\) 与 \(N_0\)，使得对一切 \(N\ge N_0\) 都有
\[
\mathcal E(N)\ge c\log N.
\]

再由定理 \ref{thm:xi-time-part9zi-fibonacci-transport-resonance-incompatibility}，
\[
F_n\,T_{F_n}(\varphi^{-1})=O(1).
\]
而由于 \(F_n\to\infty\)，前式中的对数下界给出
\[
\mathcal E(F_n)\ge c\log F_n\longrightarrow\infty.
\]

若存在某个仅依赖传输差异的控制机制，即存在函数
\[
\Phi:[0,\infty)\to[0,\infty)
\]
在有界区间上有界，并满足对全部 \(N\) 都有
\[
\mathcal E(N)\le \Phi\!\bigl(N\,T_N(\varphi^{-1})\bigr),
\]
则沿 \(N=F_n\) 应得到右端有界而左端发散的矛盾。故不存在此类控制机制。证毕。
\end{proof}

\noindent
因此，黄金分支在最佳几何均匀化尺度上的表现并不是“几何误差越小，算术记忆便越弱”的单调图景。相反，Fibonacci 长度同时暴露出两种彼此独立的极限：一方面，\(W_1\) 传输误差确实达到 \(F_n^{-1}\) 级的最优衰减；另一方面，bulk 共振梯在同一子列上却保留严格正的宏观幅度，并使累积预算至少按 \(\log N\) 增长。于是，transport regularity 与 resonance budget 之间不存在可相互代偿的统一读出准则。

\endinput
